% Generated from locale/tr/content/computability/recursive-functions/pr-relations.tex; do not hand-edit.


\olfileid[tr]{cmp}{rec}{prr}
\olsection{İlkel Özyinelemeli Bağıntılar}

\begin{defn}
Bir $R(\vec x)$ bağıntısına, karakteristik fonksiyonu
\[
\Char{R}(\vec x)=\left\{
  \begin{array}{ll}
    1 & \mbox{$R(\vec x)$ ise}\\
    0 & \mbox{aksi durumda}
  \end{array}
\right.
\]
ilkel özyinelemeliyse \emph{ilkel özyinelemeli} denir.
\end{defn}

Başka bir deyişle ilkel özyinelemeli bir $R(\vec x)$ bağıntısından söz
edildiğinde, her girdide $1$ veya $0$ döndüren ilkel özyinelemeli bir
$\Char{R}$ fonksiyonu için $\Char{R}(\vec x)=1$ biçimindeki bir bağıntı
kastedilir. Örneğin ancak ve ancak $x=0$ iken geçerli olan
$\fn{IsZero}(x)$ bağıntısı, ilkel özyinelemeyle
\begin{align*}
  \Char{\fn{IsZero}}(0) & = 1,\\
  \Char{\fn{IsZero}}(x+1) & = 0
\end{align*}
olarak tanımlanan $\Char{\fn{IsZero}}$ fonksiyonuna karşılık gelir.

Bağıntıların başka ilkel özyinelemeli fonksiyonlarla bileşkesinin
alınabileceği açıktır. Dolayısıyla aşağıdakiler de ilkel özyinelemelidir:
\begin{enumerate}
\item $\fn{IsZero}(\lvert x-y\rvert)$ ile tanımlanan eşitlik bağıntısı
  $x=y$.
\item $\fn{IsZero}(x\tsub y)$ ile tanımlanan küçük-eşit bağıntısı
  $x\leq y$.
\end{enumerate}

\begin{prop}
  İlkel özyinelemeli bağıntılar kümesi Boole işlemleri altında kapalıdır.
  Yani $P(\vec x)$ ile $Q(\vec x)$ ilkel özyinelemeliyse aşağıdakiler de
  ilkel özyinelemelidir:
  \begin{enumerate}
  \item $\lnot P(\vec x)$
  \item $P(\vec x)\land Q(\vec x)$
  \item $P(\vec x)\lor Q(\vec x)$
  \item $P(\vec x)\lif Q(\vec x)$.
  \end{enumerate}
\end{prop}

\begin{proof}
  $P(\vec x)$ ile $Q(\vec x)$ ilkel özyinelemeli, yani karakteristik
  fonksiyonları $\Char{P}$ ile $\Char{Q}$ ilkel özyinelemeli olsun.
  $\lnot P(\vec x)$ ve ötekilerin karakteristik fonksiyonlarının da ilkel
  özyinelemeli olduğunu göstermeliyiz.
  \[
    \Char{\lnot P}(\vec x)=\begin{cases}
      0 & \text{$\Char{P}(\vec x)=1$ ise}\\
      1 & \text{aksi durumda.}
    \end{cases}
  \]
  $\Char{\lnot P}(\vec x)$ fonksiyonunu
  $1\tsub\Char{P}(\vec x)$ olarak tanımlayabiliriz.
  \[
    \Char{P\land Q}(\vec x)=\begin{cases}
      1 & \text{$\Char{P}(\vec x)=\Char{Q}(\vec x)=1$ ise}\\
      0 & \text{aksi durumda.}
    \end{cases}
  \]
  $\Char{P\land Q}(\vec x)$ fonksiyonunu
  $\Char{P}(\vec x)\cdot\Char{Q}(\vec x)$ ya da
  $\fn{min}(\Char{P}(\vec x),\Char{Q}(\vec x))$ olarak tanımlayabiliriz.
  Benzer biçimde
  \begin{align*}
    \Char{P\lor Q}(\vec x) & =
      \fn{max}(\Char{P}(\vec x),\Char{Q}(\vec x)) \text{ ve}\\
    \Char{P\lif Q}(\vec x) & =
      \fn{max}(1\tsub\Char{P}(\vec x),\Char{Q}(\vec x)).
  \end{align*}
\end{proof}

\begin{prop}
  İlkel özyinelemeli bağıntılar kümesi sınırlı niceleme altında kapalıdır.
  Yani $R(\vec x,z)$ ilkel özyinelemeliyse
  \begin{align*}
    & \bforall{z<y}{R(\vec x,z)} \text{ ve}\\
    & \bexists{z<y}{R(\vec x,z)}
  \end{align*}
  bağıntıları da ilkel özyinelemelidir.
  $\bforall{z<y}{R(\vec x,z)}$, $y$'den küçük her $z$ için
  $R(\vec x,z)$ geçerli olduğunda ve ancak o zaman $\vec x$ ile $y$ için
  geçerlidir; $\bexists{z<y}{R(\vec x,z)}$ için de benzer tanım geçerlidir.
\end{prop}

\begin{proof}
  Kararlaştırma gereği $\bforall{z<0}{R(\vec x,z)}$ doğrudur; çünkü $0$'dan
  küçük hiçbir $z$ yoktur. $\bexists{z<0}{R(\vec x,z)}$ ise yanlıştır.
  Sınırlı tümel niceleyici sonlu bir çarpım ya da yinelenmiş minimum gibi
  davranır. Yani
  $P(\vec x,y)\defiff\bforall{z<y}{R(\vec x,z)}$ ise
  $\Char{P}(\vec x,y)$ şöyle tanımlanabilir:
  \begin{align*}
    \Char{P}(\vec x,0) & = 1\\
    \Char{P}(\vec x,y+1) & =
      \fn{min}(\Char{P}(\vec x,y),\Char{R}(\vec x,y)).
  \end{align*}
  Sınırlı varlıksal niceleme de benzer biçimde $\fn{max}$ kullanılarak
  tanımlanabilir. Alternatif olarak
  $\bexists{z<y}{R(\vec x,z)}\liff
  \lnot\bforall{z<y}{\lnot R(\vec x,z)}$ eşdeğerliğiyle sınırlı tümel
  nicelemeden tanımlanabilir. Örneğin
  $\bexists{x\leq y}{\dots x\dots}$ biçimindeki bir sınırlı niceleyicinin
  $\bexists{x<y+1}{\dots x\dots}$ ile eşdeğer olduğuna dikkat edin.
\end{proof}

\begin{prob}
  Üç yerli $x\equiv y\mod n$ bağıntısının ($n$ modülüne göre denklik)
  ilkel özyinelemeli olduğunu gösterin.
\end{prob}

Bir başka yararlı ilkel özyinelemeli fonksiyon da
\begin{align*}
  \fn{cond}(x,y,z) & = \begin{cases}
    y & \text{$x=0$ ise}\\
    z & \text{aksi durumda}
  \end{cases}
  \intertext{ile verilen koşullu fonksiyondur. Bu fonksiyon özyinelemeyle}
  \fn{cond}(0,y,z) & = y,\\
  \fn{cond}(x+1,y,z) & = z
\end{align*}
olarak tanımlanır. Bu fonksiyon, ilkel özyinelemeli bağıntılara göre
durumlara ayrılarak yapılan ilkel özyinelemeli fonksiyon tanımlarını
gerekçelendirmek için kullanılabilir.

\begin{prop}
$g_0(\vec x),\dots,g_m(\vec x)$ ilkel özyinelemeli fonksiyonlar ve
$R_0(\vec x),\dots,R_{m-1}(\vec x)$ ilkel özyinelemeli bağıntılar olsun.
Şu biçimde tanımlanan $f$ fonksiyonu da ilkel özyinelemelidir:
\[
  f(\vec x)=\begin{cases}
    g_0(\vec x) & \text{$R_0(\vec x)$ ise}\\
    g_1(\vec x) & \text{$R_1(\vec x)$ ve $\lnot R_0(\vec x)$ ise}\\
    \vdots & \\
    g_{m-1}(\vec x) & \text{$R_{m-1}(\vec x)$ geçerli ve önceki
      koşulların hiçbiri geçerli değilse}\\
    g_m(\vec x) & \mbox{aksi durumda.}
  \end{cases}
\]
\end{prop}

\begin{proof}
  $m=1$ iken bu, yalnızca
  \[
    f(\vec x)=\fn{cond}(\Char{\lnot R_0}(\vec x),g_0(\vec x),g_1(\vec x))
  \]
  ile tanımlanan fonksiyondur. $m>1$ için bu biçimdeki tanımların bileşkesi
  alınır.
\end{proof}
